\documentclass[11pt,a4paper]{article}
\usepackage[utf8]{inputenc}
\usepackage{authblk}
\usepackage{amsmath,amssymb,amsfonts}
\usepackage{booktabs}
\usepackage{graphicx}
\usepackage{hyperref}
\usepackage{url}
\usepackage{microtype}
\usepackage{geometry}
\geometry{margin=1in}

\title{\textbf{AgentTrace: Step-Level Hallucination Detection and Attribution in Multi-Step LLM Agent Workflows}}

\author[1]{\textbf{Somnath Chatterjee}\thanks{Research Lead}}
\author[1]{\textbf{Ayaan Khan}}
\author[1]{\textbf{Aman Hossain}}
\affil[1]{Department of Computer Science and Engineering}

\date{}

\begin{document}

\maketitle

\begin{abstract}
Large Language Model (LLM) agents are increasingly deployed for multi-step reasoning and tool-use tasks. However, these agents suffer from compounding failures, where a single step-level hallucination or reasoning error propagates through the execution trajectory, causing catastrophic task failure. Prior evaluation frameworks primarily assess agents using coarse, outcome-based success metrics or token-level hallucination flags, which fail to pinpoint exactly when and why an agent deviated from factual grounding. In this paper, we propose \textbf{AgentTrace}, a novel, cost-efficient 3-layer hybrid cascading framework designed for step-level hallucination detection, attribution, and real-time intervention. AgentTrace integrates (1) a lightweight Small Language Model (SLM) ensemble capturing context-aware semantic and contradiction signals, (2) a fine-tuned Llama-3.1-8B classifier for sequence classification and multi-class error attribution, and (3) a high-capacity Nemotron-340B judge triggered via confidence-based cascading. On the rigorous AgentHallu SOTA benchmark, AgentTrace achieves a step localization accuracy of \textbf{58.65\%}, representing a massive \textbf{+17.55\%} absolute improvement over the state-of-the-art baseline (41.1\%), while reducing latency to just \textbf{290.26 ms} on average. We demonstrate our framework through an interactive, premium-designed real-time dashboard that enables developers to visualize trajectories, identify hallucinations, and trigger active correction strategies.
\end{abstract}

\section{Introduction}
Large Language Models (LLMs) have evolved from passive text generators into autonomous agents capable of interacting with external environments through tools, APIs, and multi-step reasoning cycles (e.g., the ReAct framework~\cite{react2022}). By iteratively planning, retrieving information, executing tools, and evaluating progress, these agents solve complex tasks ranging from web navigation to scientific discovery~\cite{webarena2024, agentbench2024}. 

Despite their power, LLM agents are highly vulnerable to the \textit{error propagation cascade}. Because step $t$ in a trajectory depends directly on the history of steps $1 \dots t-1$, a single undetected hallucination—such as a fabricated tool parameter or an unjustified logical leap—rapidly contaminates subsequent reasoning. Coarse, end-to-end evaluation metrics fail to diagnose these intermediate points of failure. Conversely, token-level hallucination detection is too noisy and lacks the semantic abstraction required to identify architectural failures in tool invocation or retrieval context.

To address this challenge, the SOTA benchmark AgentHallu~\cite{agenthallu2026} established the task of \textbf{step-level hallucination detection and attribution}. Under this paradigm, a detector must analyze a multi-step trajectory and output:
\begin{enumerate}
    \item \textbf{Detection:} Whether a specific step is hallucinated.
    \item \textbf{Localization:} The exact step index $t$ where the hallucination first occurred.
    \item \textbf{Attribution:} The specific category of error out of a 6-class taxonomy (Planning, Retrieval, Reasoning, Tool-Use, Human-Interaction, or No-Hallucination).
\end{enumerate}

While high-capacity LLMs like GPT-4 or Nemotron-340B can be queried as judges to evaluate trajectories post-hoc, their high latency and cost make them impractical for real-time, in-line monitoring during agent execution. Lightweight Small Language Models (SLMs), on the other hand, are extremely fast but lack the complex reasoning capabilities to perform multi-class attribution.

In this work, we present \textbf{AgentTrace}, a highly optimized, 3-layer hybrid ensemble cascade that achieves both state-of-the-art step localization accuracy and real-time inference speed. AgentTrace implements a structured cascade that routes easy-to-detect steps using lightweight SLM features and only invokes heavy LLM judgment when step uncertainty exceeds a tight threshold. 

Our core contributions are:
\begin{itemize}
    \item We design and implement the **AgentTrace** 3-layer cascading architecture, combining an SLM feature extractor, a fine-tuned Llama-3.1-8B classifier, and a high-capacity Nemotron-340B judge.
    \item We formulate a context-aware hybrid fusion logic that integrates multiple weak signals (semantic similarity, tool claim validation, DeBERTa-based factual grounding, and sliding-window cross-step contradiction).
    \item We demonstrate that AgentTrace beats the SOTA baseline of AgentHallu by a massive \textbf{+17.55\%} absolute margin, reaching a step localization accuracy of \textbf{58.65\%} with an average latency of only \textbf{290.26 ms}.
    \item We deploy an interactive FastAPI backend and a premium Streamlit dashboard to visualize trajectories and trigger automated corrective interventions (rollback, requery, and reasoning overrides) in real-time.
\end{itemize}

\section{Related Work}
\paragraph{LLM Agent Evaluation.} Frameworks like AgentBench~\cite{agentbench2024} and WebArena~\cite{webarena2024} evaluate agents based on final task success. While highly realistic, they provide no intermediate diagnostics. AgentHallu~\cite{agenthallu2026} introduced the first step-level benchmark to measure intermediate reasoning failures, but did not propose real-time mitigation or cascading.

\paragraph{Hallucination Detection.} Traditional hallucination detection focuses on summaries or open-ended QA. These methods rely on NLI-based factual grounding against external reference texts~\cite{deberta2021} or sentence embeddings~\cite{sentencebert2019}. AgentTrace adapts these techniques to structured agent logs, validating tool inputs/outputs and tracking cross-step contradictions.

\paragraph{Cascading Inference.} Model cascading leverages a sequence of models of increasing capacity, routing queries to larger models only when smaller models express high uncertainty~\cite{llmasajudge2023}. We extend this paradigm to agent trajectory step classification, creating a cost-effective real-time defense.

\section{AgentTrace System Architecture}
The core architecture of AgentTrace consists of three cascading layers and a final decision fusion module (Figure~\ref{fig:architecture}). Let a trajectory $\mathcal{T} = \{s_1, s_2, \dots, s_N\}$ be a sequence of steps, where each step $s_t$ is defined as a tuple of the active tool, tool input, tool output, the agent's internal reasoning, and the step index $t$.

\subsection{Layer 1: SLM Ensemble (Context-Aware Signals)}
To capture high-speed semantic, syntactic, and logical indicators, we execute four parallel, lightweight SLM modules on every step $s_t$:
\begin{enumerate}
    \item \textbf{Semantic Consistency Checker:} Computes the cosine similarity between the agent's internal reasoning and the actual tool action using a lightweight Sentence-Transformer (\texttt{all-MiniLM-L6-v2}). If the similarity falls below $\theta_{\text{sim}} = 0.75$, a warning flag is raised.
    \item \textbf{Tool Claim Validator:} Extract parameter constraints and regex-based slot matches to verify that the tool output aligns with the schema required by the tool input.
    \item \textbf{Factual Grounding Engine:} Uses a fine-tuned NLI model (\texttt{nli-deberta-v3-small}) to measure the logical entailment between the step output and the external context retrieved by the agent. Steps classified as \textit{contradiction} are marked as hallucinated.
    \item \textbf{Contradiction Detector:} Measures cross-step contradiction by running a sliding NLI window comparing $s_t$ with all prior steps $s_{1 \dots t-1}$ to prevent the agent from violating its own previous claims.
\end{enumerate}

\subsection{Layer 2: Fine-Tuned Llama Classifier (QLoRA Step Attribution)}
To perform rich multi-class error attribution, we fine-tune a Llama-3.1-8B model using QLoRA~\cite{qlora2023} on step sequence classification. The step is formatted into a standardized instruction prompt containing the current trajectory state, the tool output, and the agent's reasoning. The model outputs a probability distribution over the 6-class error taxonomy:
\begin{equation}
    P(c | s_t) = \text{Softmax}(W_c \cdot h_t)
\end{equation}
where $h_t$ is the hidden representation of the step sequence and $c \in \{ \text{Planning, Retrieval, Reasoning, Tool-Use, Human-Interaction, No-Hallucination} \}$. We compute the classification confidence score as $C_{\text{Llama}} = \max_c P(c | s_t)$.

\subsection{Layer 3: Nemotron Cascading Judge (High-Uncertainty Mitigation)}
For steps where the Llama classifier is highly uncertain ($C_{\text{Llama}} < 0.70$), AgentTrace dynamically routes the trajectory to a high-capacity LLM-as-a-Judge system powered by Nemotron-340B. The Nemotron judge receives the full environment context, the execution history, and the detailed step logs, executing a multi-turn chain-of-thought prompt to verify the step and provide a structured corrective intervention.

\subsection{Signal Fusion and Threshold Cascading}
The final score for step $s_t$ is a weighted combination of the SLM ensemble signals and the Llama classifier prediction:
\begin{equation}
    S(s_t) = w_{\text{SLM}} \cdot \bar{S}_{\text{SLM}} + w_{\text{Llama}} \cdot (1 - P(\text{No-Hallucination} | s_t))
\end{equation}
where $w_{\text{SLM}} = 0.4$, $w_{\text{Llama}} = 0.6$, and $\bar{S}_{\text{SLM}}$ is the average score across the four SLM modules. A step is flagged as hallucinated if $S(s_t) \ge 0.45$.

\section{Experimental Evaluation}

\subsection{Dataset and Statistics}
We evaluate our framework on the AgentHallu trajectory benchmark alongside a high-fidelity synthetic trajectory dataset generated via OpenRouter. The dataset statistics are summarized in Table~\ref{tab:dataset_stats}.

\begin{table}[ht]
\centering
\caption{\textbf{Dataset Statistics for AgentHallu Benchmark and AgentTrace Synthetic Data}}
\label{tab:dataset_stats}
\begin{tabular}{lcc}
\toprule
\textbf{Statistic} & \textbf{AgentHallu} & \textbf{AgentTrace (Synthetic)} \\
\midrule
Total Trajectories & 500 & 200 \\
Avg Steps per Trajectory & 6.2 & 5.8 \\
Total Hallucinated Steps & 845 & 312 \\
\midrule
Planning Errors & 12\% & 6\% \\
Retrieval Errors & 25\% & 17\% \\
Reasoning Errors & 35\% & 49\% \\
Tool-Use Errors & 18\% & 20\% \\
Human-Interaction Errors & 10\% & 8\% \\
\bottomrule
\end{tabular}
\end{table}

\subsection{Main Benchmark Results}
We compare AgentTrace against the current state-of-the-art AgentHallu baseline system. The quantitative results are presented in Table~\ref{tab:main_results}.

\begin{table}[ht]
\centering
\caption{\textbf{AgentTrace vs. State-of-the-Art Baselines}}
\label{tab:main_results}
\begin{tabular}{lccc}
\toprule
\textbf{System} & \textbf{Step Localization Accuracy} & \textbf{Tool-Use Accuracy} & \textbf{FPR} \\
\midrule
AgentHallu (2026) Baseline & 41.1\% & 11.6\% & N/R \\
\textbf{AgentTrace (Ours)} & \textbf{58.65\%} & \textbf{98.0\%} & \textbf{20.3\%} \\
\bottomrule
\end{tabular}
\end{table}

AgentTrace achieves a step localization accuracy of \textbf{58.65\%}, representing a major \textbf{+17.55\%} absolute improvement. Importantly, after applying our real-time feedback loop and pipeline interventions, the agent's final tool-use execution accuracy increases to a near-perfect \textbf{98.0\%}.

\subsection{Ablation Study}
To evaluate the contribution of each SLM detector and downstream model, we conduct a systematic ablation study (Table~\ref{tab:ablation}).

\begin{table}[ht]
\centering
\caption{\textbf{Ablation Study: Impact of Detection Modules}}
\label{tab:ablation}
\begin{tabular}{lcccc}
\toprule
\textbf{Configuration} & \textbf{Loc Acc} & \textbf{Precision} & \textbf{Recall} & \textbf{F1} \\
\midrule
\textbf{Full AgentTrace} & \textbf{0.587} & \textbf{0.411} & \textbf{0.587} & \textbf{0.483} \\
w/o Contradiction Detector & 0.550 & 0.395 & 0.550 & 0.460 \\
w/o Factual Grounding & 0.520 & 0.380 & 0.520 & 0.439 \\
w/o Semantic Checker & 0.490 & 0.350 & 0.490 & 0.408 \\
w/o Tool Validator & 0.440 & 0.310 & 0.440 & 0.364 \\
\midrule
AgentHallu SOTA Baseline & 0.411 & -- & -- & -- \\
\bottomrule
\end{tabular}
\end{table}

The ablation study demonstrates that removing any single detector leads to a drop in step localization accuracy. The Tool Validator is the single most critical component, as its removal drops performance by 14.7\% absolute, indicating that verifying tool schemas and parameters is highly critical.

\subsection{Latency and Cascading Efficiency}
A key feature of AgentTrace is its computational efficiency. Running a full Nemotron-340B model on every step of a 6-step trajectory results in an average step evaluation latency of **1,850 ms**. 

By utilizing our 3-layer cascade:
\begin{itemize}
    \item \textbf{76.4\%} of steps are fully classified by the SLM ensemble and Llama-3.1-8B classifier alone.
    \item Only \textbf{23.6\%} of highly uncertain steps are escalated to Layer 3.
\end{itemize}
This cascading strategy achieves an average latency of **290.26 ms** with a P95 under **405.41 ms**, rendering the framework fully viable for real-time, interactive production use.

\section{Real-Time Intervention and Interactive UI}
To facilitate real-world adoption, AgentTrace includes a robust API backend and a production-grade interactive dashboard.

\subsection{Active Intervention Strategies}
When AgentTrace identifies a step as hallucinated, it executes one of three automated recovery strategies:
\begin{enumerate}
    \item \textbf{Tool Re-query (tool\_requery):} Re-runs the tool with corrected arguments if a parameter validation or schema mismatch is detected.
    \item \textbf{Reasoning Override (reasoning\_override):} Injects a factual premise or reasoning constraint directly into the agent's system instruction state.
    \item \textbf{Step Rollback (step\_rollback):} Reverts the agent to the last clean step ($s_{t-1}$) and forces a plan regeneration with a new random seed.
\end{enumerate}

\subsection{Interactive Streamlit Dashboard}
We implement a highly polished, responsive Streamlit dashboard (at \texttt{ui/app.py}) styled with a custom dark mode, glassmorphism containers, and live micro-animations. Developers can input a task description, trigger full-pipeline analysis on a trajectory, visualize step-by-step confidence scores, examine annotated red/green cards, and review applied intervention strategies in real-time.

\section{Conclusion and Future Work}
We introduced \textbf{AgentTrace}, a highly cost-efficient, 3-layer hybrid ensemble cascade for step-level hallucination detection and attribution. By combining fast, local SLMs with a fine-tuned Llama model and a confidence-triggered high-capacity judge, AgentTrace establishes a new SOTA on the AgentHallu benchmark (58.65\% accuracy) while maintaining sub-300 ms latencies. In future work, we plan to directly optimize agent reinforcement learning rewards using AgentTrace's step-level confidence scores.

\bibliographystyle{plain}
\bibliography{references}

\end{document}
